\subsection{Multiplicative Noise Model}
We take the following observation model:
\begin{equation}
Y_i(t) = \mu_i(t) \cdot (1 + \varepsilon_i(t))
\end{equation}
Now our observation vector is $\mathbf{Y} = [Y_1(0), Y_2(0), \ldots,
  Y_1(T-1), Y_2(T-1)]^T$. We can verify:
\begin{align}\label{eq:moments}
  \begin{split}
    \mathbb{E}[Y_i(t)] &= \mu_i(t) \cdot \mathbb{E}[1 + \varepsilon_i(t)] = \mu_i(t) \\
    \text{Var}[Y_i(t)] &= \mu_i(t)^2 \cdot \text{Var}[\varepsilon_i(t)] = \sigma^2 \mu_i(t)^2
  \end{split}
\end{align}
Then the log-likelihood is:
\begin{align}\begin{split}
\ell(\boldsymbol{\phi}) &= \sum_{t=0}^{T-1} \sum_{i=1}^{2} \left[ -\frac{1}{2}\log(2\pi\sigma^2\mu_i(t)^2) - \frac{(Y_i(t) - \mu_i(t))^2}{2\sigma^2 \mu_i(t)^2} \right] \\
&= -T\log(2\pi\sigma^2) - \sum_{t,i} \log\mu_i(t) - \frac{1}{2\sigma^2}\sum_{t,i} \frac{(Y_i(t) - \mu_i(t))^2}{\mu_i(t)^2}
\end{split}\end{align}

Taking partial with respect to $\phi_k \in \{\theta, S_1(0),
S_2(0), I_1(0), I_2(0)\}$:
\begin{align}\label{eq:partial}
  \begin{split}
  \frac{\partial \ell}{\partial \phi_k} &= -\sum_{t,i} \frac{1}{\mu_i(t)} \frac{\partial \mu_i(t)}{\partial \phi_k} + \frac{1}{2\sigma^2} \frac{\partial}{\partial \phi_k}\frac{(Y_i(t) - \mu_i(t))^2}{\mu_i(t)^2}\\
  %
  %
  %
  &= -\sum_{t,i} \frac{1}{\mu_i(t)} \frac{\partial \mu_i(t)}{\partial \phi_k} - \frac{1}{2\sigma^2}\left[\frac{2(Y_i(t) - \mu_i(t))}{\mu_i(t)^2} +\frac{2(Y_i(t) - \mu_i(t))^2}{\mu_i(t)^3}\right ] \frac{\partial \mu_i(t)}{\partial \phi_k} \\
  %
  %
  %
  &=\sum_{t,i} \frac{\partial \mu_i(t)}{\partial \phi_k} \underbrace{\left[-\frac{1}{\mu_i(t)} + \frac{1}{\sigma^2}\frac{(Y_i(t) - \mu_i(t))}{\mu_i(t)^2} + \frac{1}{\sigma^2}\frac{(Y_i(t) - \mu_i(t))^2}{\mu_i(t)^3}\right]}_{A_i(t)}
  \end{split}
\end{align}
We note that $\frac{Y_i(t) - \mu_i(t)}{\sigma\mu_i(t)} =
\frac{\epsilon_i(t)}{\sigma\mu_i(t)} \sim \mathcal{N}(0,1)$, a
standard normal random variable. Therefore:
\begin{equation}
  A_i(t) = -\frac{1}{\mu_i(t)} +\frac{Z_i(t)}{\sigma \mu_i(t)} + \frac{Z_i(t)^2}{\mu_i(t)}. 
\end{equation}
Then it is immediately evident that:
\begin{align}\begin{split}\label{eq:ind ait}
  \mathbb{E}[A_i(t)] &= 0 \\
  \mathbb{E}[A_i(t) A_j(s)] &= \mathbb{E}[A_i(t)]\mathbb{E}[A_j(s)] = 0 \text{ if } t\neq s \text{ or } i\neq j.
\end{split}\end{align}



Expanding $A_i(t)^2$:
\begin{align*}
A_i(t)^2 &= \frac{1}{\mu_i(t)^2} - \frac{2}{\sigma\mu_i(t)^2}Z_i(t) - \frac{2}{\mu_i(t)^2}Z_i(t)^2 + \frac{1}{\sigma^2\mu_i(t)^2}Z_i(t)^2 + \frac{2}{\sigma\mu_i(t)^2}Z_i(t)^3 + \frac{1}{\mu_i(t)^2}Z_i(t)^4
\end{align*}

From known Gaussian moments:
\begin{itemize}
\item $\mathbb{E}[Z_i(t)] = 0$
\item $\mathbb{E}[Z_i(t)^2] = 1$
\item $\mathbb{E}[Z_i(t)^3] = 0$ 
\item $\mathbb{E}[Z_i(t)^4] = 3$
\end{itemize}

Therefore:
\begin{align}\label{eq:var ait}
  \begin{split}
    \mathbb{E}[A_i(t)^2] &= \frac{1}{\mu_i(t)^2} - 0 - \frac{2}{\mu_i(t)^2} +
    \frac{1}{\sigma^2\mu_i(t)^2} +0 + \frac{3}{\mu_i(t)^2} \\
    %
    %
    %&= \frac{1 - 2 + 3}{\mu_i(t)^2} + \frac{1}{\sigma^2\mu_i(t)^2} \\
    %
    %
    &= \frac{2}{\mu_i(t)^2} + \frac{1}{\sigma^2\mu_i(t)^2} \\
    %
    %
    &= \frac{2\sigma^2 + 1}{\sigma^2\mu_i(t)^2}
  \end{split}
\end{align}



We can now calculate elemets of the Fisher information matrix:
\begin{align}
  \begin{split}
  J_{kl} &= \mathbb{E}\left[\frac{\partial \ell}{\partial \phi_k}
    \frac{\partial \ell}{\partial \phi_l}\right] \text{ (by the definition of
    the Fisher Information Matrix)}\\
  %
  %
  %
  &= \mathbb{E} \left [\sum_{t,i} \sum_{s,j} \frac{\partial
      \mu_i(t)}{\partial \phi_k} \frac{\partial \mu_j(s)}{\partial
      \phi_l} A_i(t) A_j(s) \right ] \text{ (eq.~\eqref{eq:partial})} \\
  %
  %
  %
  &= \sum_{t,i} \frac{\partial \mu_i(t)}{\partial
    \phi_k}\frac{\partial \mu_i(t)}{\partial \phi_l} \mathbb{E} \left
  [A_i(t)^2 \right ] \text{ (eq.~\eqref{eq:ind ait})}\\
  %
  %
  %
  &= \sum_{t,i} \frac{2\sigma^2+1}{\sigma^2\mu_i(t)^2}\frac{\partial \mu_i(t)}{\partial
    \phi_k}\frac{\partial \mu_i(t)}{\partial \phi_l}  \text{ (eq.~\eqref{eq:var ait})}\\
  \end{split}
\end{align}


In matrix notation:
\begin{equation*}
\mathbf{J} = \mathbf{G}^T \mathbf{W} \mathbf{G}
\end{equation*}
where:
\begin{itemize}
\item $\mathbf{G} = \frac{\partial \boldsymbol{\mu}}{\partial \boldsymbol{\phi}} \in \mathbb{R}^{2T \times 5}$ is the Jacobian matrix
\item $\mathbf{W} = \text{diag}(\frac{2\sigma^2 + 1}{\sigma^2\mu_1(0)^2}, \ldots, \frac{2\sigma^2 + 1}{\sigma^2\mu_2(T-1)^2}) \in \mathbb{R}^{2T \times 2T}$ is the weight matrix.
\end{itemize}
